El presente informe aborda el análisis de estabilidad y respuesta en frecuencia de circuitos realimentados propuestos en el Trabajo Práctico Nro. 3. En particular, se estudia el circuito de la figura 1 para determinar la posición de sus polos y comparar el resultado teórico con la simulación. Luego se desarrolla el diseño de la capacidad de compensación C1 para obtener un margen de fase de $45^\circ$ con $\beta=1$ y se evalúa la estabilización mediante efecto Miller y ``nulling resistor''.

A continuación, se analiza el amplificador operacional de la figura 2 utilizado como buffer de tensión, con atención a la respuesta al escalón y la obtención de un margen de fase de $60^\circ$. También se destaca la alternativa de obtención de ganancia $\text{AOL}\beta$ descrita en la figura 3, manteniendo la polarización de DC.

Cada sección incluye los resultados teóricos principales, los espacios reservados para las fórmulas y resultados numéricos, y los placeholders para las capturas de las simulaciones en LTspice.
